\scanpage{131}
% archive sheet number: 130
$V$ anneau à valuation discrète, corps résiduel $k$ parfait, corps des fractions $K$

Pour tout préschéma $X$ sur $V$, on définit $\bar X$ par Greenberg.
\[
\bar X(k')\simeq X(V_{k'})\deleted{\unsure{\longrightarrow|\longleftrightarrow}X}
\]
\margin{($k'$ alg. sur $k$, $V_{k'}$ alg. correspondante sur $V$).}

Si $k'$ est parfait, $V_{k'}$ est un anneau de valuation discrète, soit $K_{k'}=K'$ son corps des fractions, on a
\[
\bar X(\unsure{k|k'})\simeq X(V_{k'})\longrightarrow X(K')
\]

Soit \deleted{$\unsure{\varphi:}X\longrightarrow$} $Y$ un autre $V$ préschéma, \inserted{\uline{propre} \uline{sur} $V$,} et soit
\[
\varphi:\quad X_K\longrightarrow Y_K
\]
un morphisme de $K$-préschémas, nous allons en déduire une \uline{application ensembliste}
\[
\bar\varphi:\quad\bar X\longrightarrow\bar Y
\]

Soit en effet $\bar x\in\bar X$, soit $k'=\tilde{k}(\bar x)$ \deleted{une}\inserted{``la''} clôture parfaite de $k(\bar x)$, alors $\bar x$ définit un \unsure{pt|p\textsuperscript{t}} $\bar x'\in\bar X(k')$, d'où un pt \deleted{\unsure{dans} $\unsure{\varphi} + x$} $\alpha_X(\unsure{\bar x''|\bar x'})\in\bar X(V_{k'})\subset\bar X(K_{k'})$ et appliquant $\varphi$ un point
\[
\varphi(\alpha_X(\bar x'))\in\overset{\deleted{\mathord{-}}}{Y}(K_{k'})
\xleftarrow{\sim}\overset{\deleted{\mathord{-}}}{Y}(V_{k'})
\xleftarrow{\substack{\alpha_Y\\\sim}}\bar Y(k'),
\]
% slanted
\margin{Si $Y$ n'est pas propre sur $V$ et $\varphi$ \unsure{seulement} une application rationnelle, on a encore une ``application non partout définie'' $\bar X\longrightarrow\bar Y$, ayant les propriétés}
dont la localité est notre $\bar\varphi(x)$. On a donc par construction un hom
\[
k(\varphi(\bar x))\longrightarrow\widetilde{k(\bar x)}
\]
(où le $\sim$ dénote la clôture parfaite).

Pour la suite, $X,Y$ de type fini sur $K$, et \inserted{\unsure{séparés|sépr.|sérém.|sépar.}} intègres.

\uline{Proposition} 1\quad Supposons $\bar x\in\bar X$ \uline{épais} et $\varphi$ \uline{dominant \unsure{gén.\textsuperscript{t}|génér.\textsuperscript{t}} séparable} [i.e. fibre gén. séparable]. Alors\textsuperscript{\unsure{($\times$)|si}} $\bar y=\varphi(\bar x)$ \inserted{existe, et} est \uline{épais}, $k(\bar y)\subset{}k(\bar x)$, et l'application ponctuelle $\bar x_{\mathrm{réd}}\longrightarrow\bar y_{\mathrm{réd}}$ définie par cette inclusion provient d'un promorphisme d'indice fini sur un ouvert constructible dense $U\ni\bar x$, [d'image un \unsure{ens|un|ouvert} constructible de $\bar y_{\mathrm{réd}}$].
